% !TEX root = TI5.tex
\section{The complexity of $2\times 2\times 2$ tensor isomorphism over $\Z$}


The first two subsections are background on quadratic rings and Bhargava's work relating them to $2 \times 2 \times 2$ integer tensors; we prove Theorem~\ref{thm:tensorZ} in the third subsection.


\subsection{Preliminaries: Quadratic rings}

A \emph{quadratic ring} is a unital commutative ring $\rng$ of free rank $2$ as a $\Z$-module.  Such a ring is equipped with a trace: $\Tr(\alpha)$ for $\alpha \in \rng$ is the trace of the endomorphism of $\rng$ induced by $\alpha$.  The discriminant of the ring is then $\disc(\rng) = \det(\Tr(\alpha_i\alpha_j))$ where $\rng = \alpha_1 \Z + \alpha_2 \Z$.  In particular, if $\rng = \Z + \tau \Z$ where $\tau$ has minimal polynomial $a\tau^2 + b\tau + c = 0$, then $\disc(\rng) = b^2 - 4ac$.  

The discriminant is always $0, 1 \pmod{4}$ and any integer satisfying this congruence is the discriminant of a unique quadratic ring $\rng$, up to isomorphism.  In other words, discriminants parametrize quadratic rings up to isomorphism.
In particular, we can take $\rng = \Z + \tau \Z$ where $\tau$ has minimal polynomial $\tau^2 = \Delta/4$ in the case that $\Delta \equiv 0 \pmod{4}$ and $\tau^2 - \tau - \frac{\Delta-1}{4}$ in the case that $\Delta \equiv 1 \pmod{4}$.  

Quadratic rings can be listed as follows:  the ring $\Z[x]/(x^2)$ having $\disc(\rng)=0$ is the \emph{degenerate ring} (the others being \emph{non-degenerate}); the ring $(1,1)\Z + (\sqrt{\Delta},0)\Z = (1,1)\Z + \sqrt{\Delta}(\Z \oplus \Z) \subseteq \Z \oplus \Z$ of square discriminant $\disc(\rng) = \Delta \ge 1$; and the familiar \emph{number field orders} $\Z \left[ \frac{\Delta+\sqrt{\Delta}}{2} \right]$ of non-square discriminant $\disc(\rng) = \Delta$.  A ring of integers is \emph{imaginary quadratic} if $\Delta < 0$ and \emph{real quadratic} if $\Delta > 0$.

Each quadratic ring $\rng$ of non-square discriminant has an associated quadratic algebra $K = \rng \otimes_\Z \Q$ which is a number field; then $\rng$ is called a number field order.  We consider the collection of fractional ideals of $\rng$, that is, rank two $\rng$-modules in $K$.  Any fractional ideal is of the form $dI$ for $I$ an ideal of $\rng$ and $d\in K$.  The norm $N(I)$ of an ideal of $\rng$ is $|\rng/I|$, %\ynote{Sorry, this notation is unfamiliar to me. Does this mean the cardinality of $\rng/I$?}, 
and we extend this to $N(dI) = d^2 N(I)$ (which is the ratio of the covolume of the fractional ideal to that of $\rng$).  Two fractional ideals $I$ and $I'$ are considered equivalent if $I = k I'$ for some $0 \neq k \in K$.  The collection of equivalence classes of fractional ideals of $\rng$ is denoted $\Cl(\rng)$.

Bhargava states a theorem concerning the action of $\SL(2,\Z)$ instead of $\GL(2,\Z)$.  For this, he adds an orientation to a quadratic number field order $\rng$ of discriminant $\Delta$ by endowing it with a choice of $\sqrt{\Delta}$ (from among the two options).  He also adds an orientation to an ideal $I$ of $\rng$ by specifying the additional data of $\epsilon \in \{ \pm 1 \} \cong \Z/2\Z$ and performing multiplication of oriented ideals $(I,\epsilon)$ componentwise.  Scalar multiplication of oriented ideals is defined as $\alpha (I, \epsilon) = (\alpha I, \operatorname{sign}(N(\alpha))\epsilon)$ for $0 \neq \alpha \in K$.  The norm of an oriented ideal $I$ is $\epsilon N(I)$.

There is also a notion of oriented basis for $\rng$ or an ideal of $\rng$ (equivalently, for $K$).  For this, if $\tau$ is the chosen square root of $\Delta$ in the oriented ring $\rng$, then we say that $1, \tau$ is a positively oriented basis.  Any other basis $\alpha, \beta$ is defined by reference:  it is positively oriented if the change of basis matrix between $\alpha$, $\beta$ and $1, \tau$ has positive determinant; otherwise it is negatively oriented.


\subsection{Preliminaries: Bhargava's Theorem}

\begin{definition}
A triple $(I_1, I_2, I_3)$ of oriented fractional ideals of $\rng$ is considered \emph{balanced} if $$N(I_1)N(I_2)N(I_3) = 1 \text{ and } I_1I_2I_3 \subseteq \rng.$$  
Here, $N$ denotes the norm with respect to $\rng$.  Two balanced triples of ideals $(I_1, I_2, I_3)$ and $(I_1', I_2', I_3')$ for the same ring $\rng$ are considered equivalent if there exist $k_1, k_2, k_3 \in K$ such that $I_i = k_i I_i'$ for $i=1,2,3$, and orientations agree.
\end{definition}

\begin{definition}  Let $\tA \in \Z^2 \otimes \Z^2 \otimes \Z^2$ have slices $A_1$, $A_2$ with respect to one of the indices.  Then $\tA$ has \emph{discriminant} 
\[
\disc(\tA) = \disc(\det( xA_1 - yA_2)),
\]
where the right side is the discriminant of a quadratic form, namely $\disc(ax^2 + bxy + cy^2) = b^2 - 4ac$.
\end{definition}
The discriminant is preserved regardless of the direction in which we slice (the index we choose).  The tensor $\tA$ is considered \emph{non-degenerate} if it has non-zero discriminant.

Suppose $\rng = \Z + \tau \Z$ with $K = \rng \otimes_\Z \Q$ as in the previous section.  Let $\operatorname{FracI}(K)$ represent the fractional ideals of $K$.
Then, there is a map 
\[
	\rho: \Z^2 \otimes \Z^2 \otimes \Z^2 \rightarrow \operatorname{FracI}(K), 
\]
accomplished by the following process.  Take the quadratic form $Q(x,y) := ax^2 + bxy + cy^2 := -\det(xA_1 + yA_2)$.  Let $\rho(\tA)$ be the rank two $\Z$-module $\alpha \Z +  \beta \Z$ with $\rng$-module structure
\[
\tau \cdot \alpha = \frac{b+\eta}{2} \alpha + a \beta, \quad
\tau \cdot \beta = c \alpha + \frac{b-\eta}{2} \beta,
\]
where
\[
\eta = \left\{ \begin{array}{ll}
1 & \Delta \equiv 1 \pmod{4} \\
0 & \Delta \equiv 0 \pmod{4}
\end{array} \right. .
\]
The result $\rho(\tA)$ is a fractional ideal of $K$.


We consider pairs $(\rng, (I_1, I_2, I_3))$ of an oriented quadratic number field order $\rng$ and an equivalence class of balanced triples $(I_1, I_2, I_3)$ of oriented fractional ideals.  Two such pairs are \emph{isomorphic} if the rings are isomorphic and the isomorphism maps the equivalence classes to each other, preserving orientation.  

We can now state Bhargava's theorem.

\begin{theorem}[{\cite[Theorem 11]{Bhargava}}]
	\label{thm:bhargava}
	Let $\Delta$ be a non-zero non-square fundamental discriminant of a quadratic ring (necessarily a number field order).
There is a bijection between the following two sets:
\begin{enumerate}
	\item $\SL(2,\Z) \times \SL(2,\Z) \times \SL(2,\Z)$ orbits of tensors $\tA \in \Z^2 \otimes \Z^2 \otimes \Z^2$ with discriminant $\Delta$,
	\item isomorphism classes of tuples $(\rng, (I_1, I_2, I_3))$ where $\rng$ is a oriented quadratic number field order of discriminant $\Delta$, and $(I_1,I_2,I_3)$ is a balanced triple of oriented ideals.
\end{enumerate}

	The map from tensor $\tA$ to pair $(\rng, (I_1, I_2, I_3))$ is given by taking $\rng$ to be the oriented quadratic number field order with discriminant $\Delta$, taking $I_1 = \rho(\tA)$ and similarly taking $I_2$ and $I_3$ to be the analogous fractional ideals obtained by applying the definition of $\rho$ with respect to slicing along other indices.  In other words, if $\sigma$ is a $3$-cycle rotation of the indices of a tensor, then $I_i = \rho(\tA^{\sigma^i})$.  The orientation on each fractional ideal is given according to whether the basis $\alpha, \beta$ described in the definition of the map $\rho$ is positively or negatively oriented with respect to $\rng$.

	Finally, the map from pair $(\rng, (I_1,I_2,I_3))$ to tensor $\tA$ is given as follows.  Let $I_1 = \alpha_1\Z + \alpha_2\Z$, $I_2 = \beta_1\Z + \beta_2\Z$ and $I_3 = \gamma_1\Z + \gamma_2\Z$; where we have chosen the orientation of these bases to match or be different from that of $\rng = \Z + \tau \Z$ according to whether $\epsilon(I_i) = 1$ or $-1$, respectively.  Write $\alpha_i\beta_j\gamma_k = c_{ijk} + a_{ijk}\tau$, and then set the coordinates of $\tA$ to be $a_{ijk}$.
\end{theorem}

We would prefer to have a version of Bhargava's equivalence for $\GL(2,\Z)\times \GL(2,\Z) \times \GL(2,\Z)$ orbits.  Fortunately, this involves simply forgetting the orientation information. We say two pairs of the form $(\rng, (I_1, I_2, I_3))$ are \emph{unoriented isomorphic} if there is an isomorphism of rings that maps the equivalence classes to each other, but now not necessarily preserving orientation.

\begin{theorem}
	\label{thm:bhargavaGL}
	Let $\Delta$ be a non-zero non-square fundamental discriminant of a quadratic number field order.
There is a bijection between the following two sets:
\begin{enumerate}
	\item $\GL(2,\Z) \times \GL(2,\Z) \times \GL(2,\Z)$ orbits of tensors $\tA \in \Z^2 \otimes \Z^2 \otimes \Z^2$ with discriminant $\Delta$,
	\item isomorphism classes of tuples $(\rng, (I_1, I_2, I_3))$ where $\rng$ is a (unoriented) quadratic number field order of discriminant $\Delta$, and $(I_1,I_2,I_3)$ is a balanced triple of ideals.
\end{enumerate}
	The maps are as for Theorem~\ref{thm:bhargava}.
\end{theorem}

	\begin{proof}
	Observe that $\GL(2,\Z)$ is formed from $\SL(2,\Z)$ by allowing one additional transformation, namely $M = \begin{pmatrix} 0 & 1 \\ 1 & 0 \end{pmatrix}$.  An action by $M$ on the slices $A_1$ and $A_2$ swaps them, which does not change the discriminant.  This results in replacing the oriented ideal $(I_1,\epsilon)$ with $(I_1,-\epsilon)$; and fixing the ideals $(I_i,\epsilon)$ for $i=2,3$.  Thus we obtain a new pair $(\rng, (I_1', I_2', I_3'))$ which differs only in terms of orientations.  Thus the new pair is equivalent when orientations are forgotten.  

		Conversely, suppose two pairs $(R, (I_1, I_2, I_3))$ and $(R', (I_1, I_2, I_3))$ are unoriented isomorphic, i.e. when orientations are forgotten.  Let $\tA$ and $\tA'$ be the associated tensors.  Then $R \cong R'$, so without loss of generality we may assume $R = R'$.  Then by the map in Theorem~\ref{thm:bhargavaGL}, by a suitable application of $M$ in each direction to $\tA$ and $\tA'$, we can arrange that the bases of the $I_i$ are positively oriented in both cases.  So then unoriented isomorphism becomes isomorphism.
	\end{proof}

Suppose that we represent the quadratic number field order $\rng$ by its discriminant $\Delta(\rng)$, implicitly thinking of $\rng$ as having the basis $1, \tau$ as described by the minimal polynomials in the previous section.
Suppose that we then represent an ideal $I$ in terms of a basis $\alpha_1, \alpha_2$, where the $\alpha_i$ are given as $\Q$-linear combinations of $1$ and $\tau$ (i.e., an ordered pair of rationals).  This allows us to represent a tuple $(\rng, (I_1, I_2, I_3))$.
	With this representation, we observe that the maps realizing Bhargava's bijection are polynomially computable.


\subsection{Equivalence of $2 \times 2 \times 2$ TI$_\Z$ and Principality Decision}

\begin{lemma} \label{lem:tensorZ}
Suppose $\tA, \tB \in \Z^2 \otimes \Z^2 \otimes \Z^2$ have square discriminants. Then isomorphism between them (under $\GL(2,\Z) \times \GL(2,\Z) \times \GL(2,\Z)$) can be determined in polynomial time.
%Then $\tA$ is in the $\SL(2,\Z)$ orbit of either a diagonal tensor $\alpha E_{111} + \beta E_{222}$ (iff the discriminant is nonzero), or a scaled $W$-tensor $\alpha E_{111} + \alpha E_{122} + \beta E_{212}$ (iff the discriminant is zero).
\end{lemma}

A 3-hypermatrix over a field is \emph{concise} if, in each of the three directions, its slices are linearly independent; we refer to an integer hypermatrix as concise if it is concise when considered over $\Q$.

\begin{proof}
\textbf{Check the GCD.} First, the GCD of all the entries of $\tA$ is unchanged by the action by $\GL(2,\Z)^3$. If the GCD of all the entries of $\tA$ is not the same as that for $\tB$, then they are not isomorphic.

Otherwise, divide both $\tA$ and $\tB$ by the GCD. Because this division was action by a scalar, it commutes with the action of $\GL(2,\Z)^3$, so the original $\tA$ and $\tB$ were isomorphic iff the rescaled versions are. We may thus now assume that $\tA$ (resp. $\tB$) has the property that the GCD of its entries is 1.

\textbf{The non-concise case.} Next, let $A_1, A_2$ be the frontal slices of $\tA$. If $A_2$ is a $\Q$-scalar multiple of $A_1$, then we can take the GCD, replace $A_1$ by some scalar multiple of itself and $A_2$ by zero. In this case, we check that $\tB$ has the same structure, and thus reduce to the case of the $\GL(2,\Z) \times \GL(2,\Z)$ equivalence of two $2 \times 2$ matrices under the left-right action. This is solved by putting both matrices into Smith Normal Form and checking that the two have equal SNFs.

Thus, we now assume that $A_1$ and $A_2$ are not $\Q$-multiples of one another (and analogously for the slices in the other two directions).

\textbf{The concise case.} Since the discriminant $D$ is a perfect square, the form $\det(x A_1 - y A_2) = ax^2 + bxy + cy^2$ has two integer roots up to scale (which need not be distinct, namely if $D=0$), given by $(x,y) = (-b + \sqrt{D}, 2a)$ and $(x,y) = (-b - \sqrt{D}, 2a)$. Now, let $\varepsilon$ be the GCD of $-b + \sqrt{D}$ and $2a$, $\alpha = (-b + \sqrt{D})/\varepsilon$, $\beta = -2a / \varepsilon$. Since $\alpha$ and $\beta$ are coprime, there exist $\gamma,\delta \in \Z$ such that $\alpha \delta - \beta \gamma = 1$. Define $A_1', A_2'$ by
\[
\begin{pmatrix}
A_1' \\
A_2'
\end{pmatrix}
=
\begin{pmatrix}
\alpha & \beta \\
\gamma & \delta
\end{pmatrix}
\begin{pmatrix}
A_1 \\
A_2
\end{pmatrix}.
\]
Then the tensor $\tA'$ whose frontal slices are $A_1',A_2'$ differs from $\tA$ by an element of $\SL(2,\Z)$. 

Now, by construction we have $\det(A_1') = 0$. Use the action of $\GL(2,\Z) \times \GL(2,\Z)$ by left and right multiplication to put $A_1'$ into Smith Normal Form $A_1''$, while simultaneously transforming $A_2'$ into some matrix $A_2''$. Since $A_1'$ has rank at most 1, its Smith Normal Form is 
\[
A_1'' = \begin{pmatrix} e & 0 \\ 0 & 0 \end{pmatrix}
\]
for some integer $e$.
Let $\tA''$ be the tensor with frontal slices $A_1'', A_2''$, which is again isomorphic to $\tA$. We now have
\[
\det(x A_1'' - y A_2'') = e (A_2'')_{22} xy + \det(A_2'') y^2 = y(e (A_2'')_{22}  x + \det(A_2'') y),
\]
Since the tensor is concise, we have $e \neq 0$, and thus $D = 0 \Leftrightarrow (A_2'')_{22} = 0$. 


\textbf{Concise, discriminant zero.} Suppose that $D = 0 = (A_2'')_{22}$.  Then $\det(x A_1'' - y A_2'')$ has a double root at $(1,0)$, so $A_1''$ and its scalar multiples are the only rank-1 matrices in $Span_\Z(A_1'', A_2'')$. Next, the stabilizer of $A_1''$ under the left-right action of $\GL(2,\Z) \times \GL(2,\Z)$ consists of those pairs $(P,Q) \in \GL(2,\Z)^2$ such that $P,Q$ are both upper triangular (where the action is $(P,Q) \cdot A_1'' = P A_1'' Q^T$). Thus we may use this stabilizer to fix $A_1''$ while adding multiples of the (1,2) and (2,1) entries to the (1,1) entry of $A_2''$. We may also add multiples of $A_1''$ to $A_2''$. Thus, we arrive at the form
\[
A_1'' = \begin{pmatrix} e & 0 \\ 0 & 0 \end{pmatrix} \qquad A_2'' = \begin{pmatrix} f & g \\ h & 0 \end{pmatrix},
\]
where $f$ is reduced modulo the GCD of $g$, $h$, and $e$, and at least one of $g$ or $h$ is nonzero (since otherwise $A_1''$ would be a $\Q$-multiple of $A_2''$, which we handled above).

Now, attempt to do the same procedure to $\tB$. And suppose we arrive at 
\[
B_1'' = \begin{pmatrix} e' & 0 \\ 0 & 0 \end{pmatrix} \qquad B_2'' = \begin{pmatrix} f' & g' \\ h' & 0 \end{pmatrix}.
\]
If any part of this procedure does not result in this form, then $\tA$ and $\tB$ are not isomorphic. Since $A_1''$ and its multiples are the only rank-1 matrices in $Span_\Z(A_1'', A_2'')$, and the same is true of $B_1''$, if $e \neq \pm e'$, then $\tA$ and $\tB$ cannot be isomorphic. Otherwise, if needed, multiply by an element of $\GL(2,\Z)$ to fix the sign so that $e = e'$. 

At this point, if $(P,Q,R)$ is an isomorphism from $\tA''$ to $\tB''$, $R$ must fix the first slice (hence be lower triangular), and thus the action of $(P,Q)$ must also fix the first slice, so $P$ and $Q$ must both be upper triangular. Under the action of this subgroup, the only things we can do are to add multiples of $e,g,h$ to $f$. If the pair $(g,h)$ is not equal to one of $(\pm g', \pm h')$, then $\tA$ cannot be isomorphic to $\tB$. Otherwise, multiply the second row and second column by the necessary signs to make $g = g'$ and $h = h'$. Finally, since $f$ and $f'$ are both reduced to their minimal representatives modulo the ideal $\langle e,g,h \rangle = \langle e', g', h' \rangle$ (equivalently, modulo the GCD of $e$, $g$, and $h$), we have that $\tA \cong \tB$ if and only if $f = \pm f'$. If the signs are not equal, note that we can multiply the second slice, the second row, and the second column all by $-1$, and this has the effect of only multiplying the $f$ entry by $-1$. 

\textbf{Concise, non-zero discriminant.} On the other hand, suppose that both $D$ and $(A_2'')_{22}$ are nonzero. Now there are two $\Q$-linearly independent matrices of rank 1 in $Span_\Z(A_1, A_2)$. We may either replace $A_1$ by such a matrix $A_1'$ and get a resulting $A_2'$ as above, or vice versa. We handle the first case here, the other case being the same just swapping the roles of $A_1$ and $A_2$, \emph{mutatis mutandis}. 

Thus, we have arrived at a form
\[
A_1'' = \begin{pmatrix} e & 0 \\ 0 & 0 \end{pmatrix} \qquad A_2'' = \begin{pmatrix} f & g \\ h & k \end{pmatrix},
\]
with $k \neq 0$, and similarly
\[
B_1'' = \begin{pmatrix} e' & 0 \\ 0 & 0 \end{pmatrix} \qquad B_2'' = \begin{pmatrix} f' & g' \\ h' & k' \end{pmatrix}.
\]
(If $\tB$ cannot be put in this form, then the two are not isomorphic.) As above, to preserve the slice $A_1''$, we cannot add any multiple of $A_2''$ to $A_1''$, and we can only apply the left-right action of $(P,Q)$, where both $P$ and $Q$ are upper triangular. Since we cannot alter $A_1''$, if $e \neq \pm e'$, then we try putting $\tB$ into this form using the other rank-1 matrix in the span of its slices. And if neither works, then the two are not isomorphic. 

So we may now assume $e = e'$ (by multiplying by a sign if needed). Then the two are isomorphic, by an isomorphism that sends $A_1''$ to $B_1''$, if and only if there is an integer solution to: % \ynote{an integer solution?} to:
\[
\begin{pmatrix} 1 & x \\ 0 & 1 \end{pmatrix}
\begin{pmatrix} f + ze & g \\ h & k \end{pmatrix}
=
\begin{pmatrix} f' & g' \\ h' & k' \end{pmatrix}
\begin{pmatrix} 1 & 0 \\ y & 1 \end{pmatrix}.
\]
This is a set of integer linear equations in $x$, $y$, and $z$ and thus can be efficiently solved.

(In more detail, one way to see the solution directly, instead of using a generic algorithm for solving linear Diophantine equations, is as follows. The action cannot change $k$ and $k'$, so we need $k = \pm k'$, and again we may multiply by signs to achieve $k = k'$. Using $x$, we can take $g$ modulo $k$, so we must have $g \equiv g' \pmod{k}$. Using $y$, we can take $h'$ modulo $k$, so we must have $h \equiv h' \pmod{k}$. And after we use $x$ and $y$ in this way, whatever the new values in the positions of $f$ and $f'$ are, they must be equivalent modulo $e$.)
\end{proof}




\begin{proof}[Proof of Theorem~\ref{thm:tensorZ}]
Let $\tA$ and $\tB$ be two elements of $\Z^2 \otimes \Z^2 \otimes \Z^2$.  
Compute $\disc(\tA)$ and $\disc(\tB)$.  If they are not equal, return false.  
If they are equal and square, solve Tensor Isomorphism using Lemma~\ref{lem:tensorZ}.

	It remains to consider the case that $\disc(\tA) = \disc(\tB)$ is non-square.
In this case, we will reduce \algprobm{Tensor Isomorphism} to \algprobm{Principality Decision}.
	Write $\rng$ for the associated quadratic ring of discriminant $\disc(\tA)$.
	We will determine instead if the associated tuples $(\rng, (I_1,I_2,I_3))$ are unoriented isomorphic, which suffices by Theorem~\ref{thm:bhargavaGL}.
Compute $I_1(\tA)$ and $I_1(\tB)$.  These are ideals of $\rng$.  

	Next, we show that to determine whether $I_1(\tA)$ and $I_1(\tB)$ are equivalent ideals of $\rng$ is to check an ideal for principality.
	First, recall that the \emph{order} of an ideal $I$ is the ring $\operatorname{ord}(I) := \{ \alpha \in K : \alpha I \subseteq I \}$.  This can be computed polynomially by a linear algebra algorithm, or by reducing to computing the GCD of the coefficients of the associated quadratic forms \cite[Section 7.A]{Cox}.  If the orders of $I_1(\tA)$ and $I_1(\tB)$ are not the same, the ideals are not equivalent.  If the orders are the same, then call the order $\rng'$ and observe that $\rng \subseteq \rng'$.  In this case the two ideals are invertible in $\rng'$.   
	To determine if they are equal, it suffices to test $I_i(\tA) I_i(\tB)^{-1}$ for principality.

	Assuming we can test for principality, if the forms are not equal, return false.  If they are equal, proceed to test $I_2(\tA)$ and $I_2(\tB)$.  If not equivalent, return false.  If they are equal, proceed to test $I_3(\tA)$ and $I_3(\tB)$.  If equivalent, return true; otherwise, return false.  

	As far as principality testing, if $\disc(\tA) < 0$, then invertible ideals correspond to positive definite primitive integral binary quadratic forms via a polynomially computable isomorphism, and the classical reduction theory of positive definite primitive integral binary quadratic forms can check efficiently whether an ideal is principal \cite[Section 5.3.1]{Cohen}.

	If $\disc(\tA) > 0$, we are in the real quadratic case, and we rely on the Principality Decision Problem to test if $I_i(\tA)$ and $I_i(\tB)$ are equivalent.  
	Since all cases except for this last one can be solved in polynomial time, this gives a polynomial-time reduction to Principality Decision for ideals in real quadratic fields.


Conversely, we reduce \algprobm{Principality Decision} in real quadratic fields to \algprobm{Tensor Isomorphism}.  Let $I$ be an ideal of a real quadratic number field order $\rng$, given by an integral basis.
We will rephrase the principality problem for $I$ as a tensor problem.  First, compute the inverse of $I$, which is $\overline I/N(I)$, which is easily computable from a basis.  Observe that the triple $(\rng, (I, I^{-1},R))$ is balanced, since $I I^{-1} R = R$.  Then we form $\tA$ from the triple $(\rng, (I,I^{-1},R))$ according to Theorem~\ref{thm:bhargavaGL}, which takes polynomial time.  We also form $\tB$ from the trivial ideal $(\rng, (\rng, \rng, \rng))$ (which is also balanced).  Then, $I$ is principal if and only if $\tA$ and $\tB$ are isomorphic.

	Finally, with the above reduction, the statement about $\cc{BQP}$ then follows immediately from \cite[Theorem 4.1]{Hallgren}.
\end{proof}


